\documentclass[11pt,a4paper]{article}

\usepackage[a4paper,margin=20mm]{geometry}
\usepackage{amsmath,amssymb}
\usepackage{booktabs,array}
\usepackage[hidelinks]{hyperref}

\setlength{\parskip}{6pt plus 2pt minus 1pt}
\setlength{\parindent}{0pt}

\newcommand{\optsfourA}[4]{%
  \begin{tabbing}
    \hspace{0.48\linewidth}\=\kill
    (A) #1 \> (B) #2 \\[2pt]
    (C) #3 \> (D) #4
  \end{tabbing}%
}

\newcommand{\optslistA}[4]{%
  \begin{itemize}
    \setlength{\itemsep}{2pt}\setlength{\parskip}{0pt}
    \item[(A)] #1
    \item[(B)] #2
    \item[(C)] #3
    \item[(D)] #4
  \end{itemize}%
}

\newcommand{\ans}[1]{\textbf{Answer:} #1\par\medskip}

\begin{document}

% ============================================================================
% GATE 2025 -- GENERAL APTITUDE (GA) & MATHEMATICS (MA)
% ============================================================================

\begin{center}
  {\LARGE\bfseries GATE 2025}\\[4pt]
  {\large\bfseries General Aptitude (GA) \& Mathematics (MA)}\\[2pt]
  {\small\itshape Official Final Answer Key attached to every question}
\end{center}
\medskip

\section*{General Aptitude (GA)}

\subsection*{Q.1 -- Q.5 Multiple Choice Question (MCQ), carry ONE mark each (for each wrong answer: $-1/3$).}

\textbf{Q.1} Ravi had \underline{\hspace{1.5cm}} younger brother who taught at \underline{\hspace{1.5cm}} university. He was widely regarded as \underline{\hspace{1.5cm}} honorable man.\\
Select the option with the correct sequence of articles to fill in the blanks.
\optsfourA{a; a; an}{the; an; a}{a; an; a}{an; an; a}
\ans{(A)}

\textbf{Q.2} The CEO's decision to downsize the workforce was considered myopic because it sacrificed long-term stability to accommodate short-term gains.\\
Select the most appropriate option that can replace the word ``myopic'' without changing the meaning of the sentence.
\optsfourA{visionary}{shortsighted}{progressive}{innovative}
\ans{(B)}

\textbf{Q.3} The average marks obtained by a class in an examination were calculated as $30.8$. However, while checking the marks entered, the teacher found that the marks of one student were entered incorrectly as $24$ instead of $42$. After correcting the marks, the average becomes $31.4$. How many students does the class have?
\optsfourA{$25$}{$28$}{$30$}{$32$}
\ans{(C)}

\textbf{Q.4} Consider the relationships among P, Q, R, S, and T:
\begin{itemize}
  \setlength{\itemsep}{1pt}\setlength{\parskip}{0pt}
  \item P is the brother of Q.
  \item S is the daughter of Q.
  \item T is the sister of S.
  \item R is the mother of Q.
\end{itemize}
The following statements are made based on the relationships given above.
\begin{itemize}
  \setlength{\itemsep}{1pt}\setlength{\parskip}{0pt}
  \item[(1)] R is the grandmother of S.
  \item[(2)] P is the uncle of S and T.
  \item[(3)] R has only one son.
  \item[(4)] Q has only one daughter.
\end{itemize}
Which one of the following options is correct?
\optsfourA{Both (1) and (2) are true.}{Both (1) and (3) are true.}{Only (3) is true.}{Only (4) is true.}
\ans{(A)}

\textbf{Q.5} According to the map shown in the figure, which one of the following statements is correct?\\
\emph{[Map layout (North up, East right): 1st Main Road runs North--South and intersects 5th Cross Road running West--East. Northwest quadrant (west of 1st Main Road, north of 5th Cross Road): Library on top, Physics Lab directly south of Library. Northeast quadrant: Canteen on top (due east of Library), Hospital directly south of Canteen (due east of Physics Lab). Southwest quadrant: Hostels on top, Classrooms directly south of Hostels. Southeast quadrant: Chemistry Lab (directly south of Hospital and southeast of Physics Lab).]}
\optslistA{The library is located to the northwest of the canteen.}{The hospital is located to the east of the chemistry lab.}{The chemistry lab is to the southeast of physics lab.}{The classrooms and canteen are next to each other.}
\ans{(C)}

\subsection*{Q.6 -- Q.10 Multiple Choice Question (MCQ), carry TWO marks each (for each wrong answer: $-2/3$).}

\textbf{Q.6} ``I put the brown paper in my pocket along with the chalks, and possibly other things. I suppose every one must have reflected how primeval and how poetical are the things that one carries in one's pocket: the pocket-knife, for instance the type of all human tools, the infant of the sword. Once I planned to write a book of poems entirely about the things in my pocket. But I found it would be too long: and the age of the great epics is past.'' \hfill (From G.K. Chesterton's ``A Piece of Chalk'')\\
Based only on the information provided in the above passage, which one of the following statements is true?
\optslistA{The author of the passage carries a mirror in his pocket to reflect upon things.}{The author of the passage had decided to write a poem on epics.}{The pocket-knife is described as the infant of the sword.}{Epics are described as too inconvenient to write.}
\ans{(C)}

\textbf{Q.7} In the diagram, the lines $QR$ and $ST$ are parallel to each other. The shortest distance between these two lines is half the shortest distance between the point $P$ and line $QR$. What is the ratio of the area of the triangle $PST$ to the area of the trapezium $SQRT$?\\
\emph{[Diagram note: Triangle $PQR$ with apex $P$ at the top and horizontal base $QR$ at the bottom; segment $ST$ is parallel to $QR$ with $S$ on side $PQ$ and $T$ on side $PR$, halfway between $P$ and $QR$.]}
\optsfourA{$\dfrac{1}{3}$}{$\dfrac{1}{4}$}{$\dfrac{2}{5}$}{$\dfrac{1}{2}$}
\ans{(A)}

\textbf{Q.8} A fair six-faced dice, with the faces labelled `1', `2', `3', `4', `5', and `6', is rolled thrice. What is the probability of rolling `6' exactly once?
\optsfourA{$\dfrac{75}{216}$}{$\dfrac{1}{6}$}{$\dfrac{1}{18}$}{$\dfrac{25}{216}$}
\ans{(A)}

\textbf{Q.9} A square paper, shown in figure (I), is folded along the dotted lines as shown in the figures (II) and (III). Then a few cuts are made as shown in figure (IV). Which one of the following patterns will be obtained when the paper is unfolded?\\
\emph{[Diagram note: Square (I) is folded diagonally along the top-left to bottom-right diagonal (II) to form an upper-right triangle, then folded along the perpendicular bisector from the top-right corner to the diagonal midpoint (III) to form a smaller isosceles right triangle (IV) with its right-angle corner at the bottom (corresponding to the original top-right/bottom-left corners) and hypotenuse along the top. Three cuts are made in (IV): a small triangle pointing up near the center of the top edge, a small square/rectangle on the left slant edge, and a small rectangle perpendicular to the right slant edge. In option (A), unfolding produces four outward-pointing small triangles at the midpoints of all four outer edges, two small squares near the top-left and bottom-right corners, and two L-shaped cuts near the top-right and bottom-left corners; (B) replaces the corner squares with rotated diamonds; (C) points the top/bottom triangles inward; (D) places a single square in the center instead of at the corners.]}
\optslistA{Pattern (A): four outward-pointing edge-midpoint triangles, two squares at top-left/bottom-right, two L-cuts at top-right/bottom-left}{Pattern (B): four outward-pointing triangles, two rotated diamonds at top-left/bottom-right, two L-cuts}{Pattern (C): inward-pointing top/bottom triangles with two corner squares and two L-cuts}{Pattern (D): four outward-pointing triangles, one center square, two L-cuts}
\ans{(A)}

\textbf{Q.10} A shop has $4$ distinct flavors of ice-cream. One can purchase any number of scoops of any flavor. The order in which the scoops are purchased is inconsequential. If one wants to purchase $3$ scoops of ice-cream, in how many ways can one make that purchase?
\optsfourA{$4$}{$20$}{$24$}{$48$}
\ans{(B)}

\section*{Mathematics (MA)}

\subsection*{Q.11 -- Q.22 Multiple Choice Question (MCQ), carry ONE mark each (for each wrong answer: $-1/3$).}

\textbf{Q.11} Let
\[
  S = \left\{w = \begin{bmatrix} w_1 \\ w_2 \\ w_3 \end{bmatrix} \in \mathbb{R}^3 : \begin{bmatrix} 1 \\ -3 \\ 2 \end{bmatrix} \begin{bmatrix} w_1 & w_2 & w_3 \end{bmatrix} \text{ is diagonalizable and } \|w\| = 1\right\},
\]
where $\|w\| = (w_1^2 + w_2^2 + w_3^2)^{1/2}$. Then, which one of the following is TRUE?
\optslistA{$S$ is compact and connected}{$S$ is neither compact nor connected}{$S$ is compact but not connected}{$S$ is connected but not compact}
\ans{(B)}

\textbf{Q.12} Given that the Laplace transforms of $J_0(x)$, $J_0'(x)$ and $J_0''(x)$ exist, where $J_0(x)$ is the Bessel function. Let $Y = Y(s)$ be the Laplace transform of the Bessel function $J_0(x)$. Then, which one of the following is TRUE?
\optsfourA{$\dfrac{dY}{ds} + \dfrac{2sY}{s^2 + 1} = 0, \; s > 0$}{$\dfrac{dY}{ds} - \dfrac{2sY}{s^2 + 1} = 0, \; s > 0$}{$\dfrac{dY}{ds} - \dfrac{sY}{s^2 + 1} = 0, \; s > 0$}{$\dfrac{dY}{ds} + \dfrac{sY}{s^2 + 1} = 0, \; s > 0$}
\ans{(D)}

\textbf{Q.13} To find a real root of the equation $x^3 + 4x^2 - 10 = 0$ in the interval $(1, \frac{3}{2})$ by using the fixed-point iteration scheme, consider the following two statements:\\
S1: The iteration scheme $x_{k+1} = \sqrt{\dfrac{10}{4 + x_k}}$, $k = 0, 1, 2, \ldots$, converges for any initial guess $x_0 \in (1, \frac{3}{2})$.\\
S2: The iteration scheme $x_{k+1} = \dfrac{1}{2}\sqrt{10 - x_k^3}$, $k = 0, 1, 2, \ldots$, diverges for some initial guess $x_0 \in (1, \frac{3}{2})$.\\
Then, which one of the following is correct?
\optslistA{S1 is TRUE and S2 is FALSE}{S2 is TRUE and S1 is FALSE}{both S1 and S2 are TRUE}{neither S1 nor S2 is TRUE}
\ans{(A)}

\textbf{Q.14} For the linear programming problem:
\begin{align*}
  \text{Maximize} \quad & Z = 2x_1 + 4x_2 + 4x_3 - 3x_4 \\
  \text{Subject to} \quad & \alpha x_1 + x_2 + x_3 = 4, \\
  & x_1 + \beta x_2 + x_4 = 8, \\
  & x_1, x_2, x_3, x_4 \ge 0,
\end{align*}
consider the following two statements:\\
S1: If $\alpha = 2$ and $\beta = 1$, then $(x_1, x_2)^T$ forms an optimal basis.\\
S2: If $\alpha = 1$ and $\beta = 4$, then $(x_3, x_2)^T$ forms an optimal basis.\\
Then, which one of the following is correct?
\optslistA{S1 is TRUE and S2 is FALSE}{S2 is TRUE and S1 is FALSE}{both S1 and S2 are TRUE}{neither S1 nor S2 is TRUE}
\ans{(B)}

\textbf{Q.15} Consider the following subsets of the Euclidean space $\mathbb{R}^4$:
\begin{align*}
  S &= \{(x_1, x_2, x_3, x_4) \in \mathbb{R}^4 : x_1^2 + x_2^2 + x_3^2 - x_4^2 = 0\}, \\
  T &= \{(x_1, x_2, x_3, x_4) \in \mathbb{R}^4 : x_1^2 + x_2^2 + x_3^2 - x_4^2 = 1\}, \quad \text{and} \\
  U &= \{(x_1, x_2, x_3, x_4) \in \mathbb{R}^4 : x_1^2 + x_2^2 + x_3^2 - x_4^2 = -1\}.
\end{align*}
Then, which one of the following is TRUE?
\optslistA{$S$ is connected, but $T$ and $U$ are not connected}{$T$ and $U$ are connected, but $S$ is not connected}{$S$ and $U$ are connected, but $T$ is not connected}{$S$ and $T$ are connected, but $U$ is not connected}
\ans{(D)}

\textbf{Q.16} Consider the system of ordinary differential equations
\[
  \frac{dX}{dt} = MX,
\]
where $M$ is a $6 \times 6$ skew-symmetric matrix with entries in $\mathbb{R}$. Then, for this system, the origin is a stable critical point for
\optslistA{any such matrix $M$}{only such matrices $M$ whose rank is $2$}{only such matrices $M$ whose rank is $4$}{only such matrices $M$ whose rank is $6$}
\ans{(A)}

\textbf{Q.17} Let $X = \{f \in C[0, 1] : f(0) = 0 = f(1)\}$ with the norm $\|f\|_\infty = \sup_{0 \le t \le 1} |f(t)|$, where $C[0, 1]$ is the space of all real-valued continuous functions on $[0, 1]$. Let $Y = C[0, 1]$ with the norm $\|f\|_2 = (\int_0^1 |f(t)|^2\,dt)^{1/2}$. Let $U_X$ and $U_Y$ be the closed unit balls in $X$ and $Y$ centred at the origin, respectively. Consider $T : X \to \mathbb{R}$ and $S : Y \to \mathbb{R}$ given by
\[
  Tf = \int_0^1 f(t)\,dt \quad \text{and} \quad Sf = \int_0^1 f(t)\,dt.
\]
Consider the following statements:\\
S1: $\sup_{f \in U_X} |Tf|$ is attained at a point of $U_X$.\\
S2: $\sup_{f \in U_Y} |Sf|$ is attained at a point of $U_Y$.\\
Then, which one of the following is correct?
\optslistA{S1 is TRUE and S2 is FALSE}{S2 is TRUE and S1 is FALSE}{both S1 and S2 are TRUE}{neither S1 nor S2 is TRUE}
\ans{(B)}

\textbf{Q.18} Let $g(x, y) = f(x, y) e^{2x+3y}$ be defined in $\mathbb{R}^2$, where $f(x, y)$ is a continuously differentiable non-zero homogeneous function of degree $4$. Then,
\[
  x \frac{\partial g}{\partial x} + y \frac{\partial g}{\partial y} = 0
\]
holds for
\optslistA{all points $(x, y)$ in $\mathbb{R}^2$}{all points $(x, y)$ on the line given by $2x + 3y + 4 = 0$}{all points $(x, y)$ in the region of $\mathbb{R}^2$ except on the line given by $2x + 3y + 4 = 0$}{all points $(x, y)$ on the line given by $2x + 3y = 0$}
\ans{(B)}

\textbf{Q.19} The partial differential equation
\[
  (1 + x^2) \frac{\partial^2 u}{\partial x^2} + 2x(1 - y^2) \frac{\partial^2 u}{\partial y \partial x} + (1 - y^2) \frac{\partial^2 u}{\partial y^2} + x \frac{\partial u}{\partial x} + (1 - y^2) \frac{\partial u}{\partial y} = 0
\]
is
\optslistA{elliptic in the region $\{(x, y) \in \mathbb{R}^2 : |y| \le 1\}$}{hyperbolic in the region $\{(x, y) \in \mathbb{R}^2 : |y| > 1\}$}{elliptic in the region $\{(x, y) \in \mathbb{R}^2 : |y| > 1\}$}{hyperbolic in the region $\{(x, y) \in \mathbb{R}^2 : |y| < 1\}$}
\ans{(B)}

\textbf{Q.20} Let $u(x, t)$ be the solution of the following initial-boundary value problem
\[
  \frac{\partial u}{\partial t} - \frac{\partial^2 u}{\partial x^2} = 0, \quad x \in (0, \pi), \; t > 0, \qquad u(0, t) = u(\pi, t) = 0, \quad u(x, 0) = \sin 4x \cos 3x.
\]
Then, for each $t > 0$, the value of $u\!\left(\frac{\pi}{4}, t\right)$ is
\optsfourA{$\dfrac{e^{-49t}}{2\sqrt{2}}(e^{48t} - 1)$}{$\dfrac{e^{-49t}}{2\sqrt{2}}(1 - e^{48t})$}{$\dfrac{e^{-49t}}{2\sqrt{2}}(1 + e^{48t})$}{$\dfrac{e^{-49t}}{4\sqrt{2}}(1 - e^{48t})$}
\ans{(A)}

\textbf{Q.21} Consider the function $F : \mathbb{R}^2 \to \mathbb{R}^2$ given by
\[
  F(x, y) = (x^3 - 3xy^2 - 3x, \; 3x^2 y - y^3 - 3y).
\]
Then, for the function $F$, the inverse function theorem is
\optslistA{applicable at all points of $\mathbb{R}^2$}{not applicable at exactly one point of $\mathbb{R}^2$}{not applicable at exactly two points of $\mathbb{R}^2$}{not applicable at exactly three points of $\mathbb{R}^2$}
\ans{(C)}

\textbf{Q.22} Let the functions $f : \mathbb{R}^2 \to \mathbb{R}$ and $g : \mathbb{R}^2 \to \mathbb{R}$ be given by
\[
  f(x_1, x_2) = x_1^2 + x_2^2 - 2x_1 x_2 \quad \text{and} \quad g(x_1, x_2) = 2x_1^2 + 2x_2^2 - x_1 x_2.
\]
Consider the following statements:\\
S1: For every compact subset $K$ of $\mathbb{R}$, $f^{-1}(K)$ is compact.\\
S2: For every compact subset $K$ of $\mathbb{R}$, $g^{-1}(K)$ is compact.\\
Then, which one of the following is correct?
\optslistA{S1 is TRUE and S2 is FALSE}{S2 is TRUE and S1 is FALSE}{both S1 and S2 are TRUE}{neither S1 nor S2 is TRUE}
\ans{(B)}

\subsection*{Q.23 -- Q.26 Multiple Select Question (MSQ), carry ONE mark each (no negative marking).}

\textbf{Q.23 (MSQ)} Let $p_A(x)$ denote the characteristic polynomial of a square matrix $A$. Then, for which of the following invertible matrices $M$, the polynomial $p_M(x) - p_{M^{-1}}(x)$ is constant?
\optsfourA{$M = \begin{bmatrix} 5 & 7 \\ 2 & 3 \end{bmatrix}$}{$M = \begin{bmatrix} 3 & 1 \\ 4 & 2 \end{bmatrix}$}{$M = \begin{bmatrix} 1 & 3 \\ 2 & -1 \end{bmatrix}$}{$M = \begin{bmatrix} 5 & -8 \\ 2 & -3 \end{bmatrix}$}
\ans{(A), (C), (D) (MSQ)}

\textbf{Q.24 (MSQ)} Consider the balanced transportation problem with three sources $S_1, S_2, S_3$, and four destinations $D_1, D_2, D_3, D_4$, for minimizing the total transportation cost whose cost matrix is as follows:
\begin{center}
\begin{tabular}{cccccc}
\toprule
& $D_1$ & $D_2$ & $D_3$ & $D_4$ & \textbf{Supply} \\
\midrule
$S_1$ & $2$ & $6$ & $20$ & $11$ & $\alpha + 10$ \\
$S_2$ & $12$ & $7$ & $4$ & $10$ & $\alpha + \lambda + 10$ \\
$S_3$ & $8$ & $14$ & $16$ & $11$ & $5$ \\
\midrule
\textbf{Demand} & $\alpha + 5$ & $10$ & $\lambda + 5$ & $\alpha + \lambda$ & \\
\bottomrule
\end{tabular}
\end{center}
where $\alpha, \lambda > 0$. If the associated cost to the starting basic feasible solution obtained by using the North-West corner rule is $290$, then which of the following is/are correct?
\optslistA{$\alpha^2 + \lambda^2 = 100$}{$\alpha^2 + \alpha\lambda = 150$}{The optimal cost of the transportation problem is $260$}{The optimal cost of the transportation problem is $290$}
\ans{(B), (D) (MSQ)}

\textbf{Q.25 (MSQ)} Consider the following regions:
\begin{align*}
  S_1 &= \{(x_1, x_2) \in \mathbb{R}^2 : 2x_1 + x_2 \le 4, \; x_1 + 2x_2 \le 5, \; x_1, x_2 \ge 0\}, \\
  S_2 &= \{(x_1, x_2) \in \mathbb{R}^2 : 2x_1 - x_2 \le 5, \; x_1 + 2x_2 \le 5, \; x_1, x_2 \ge 0\}.
\end{align*}
Then, which of the following is/are TRUE?
\optslistA{The maximum value of $x_1 + x_2$ is $3$ on the region $S_2$}{The maximum value of $x_1 + x_2$ is $5$ on the region $S_2 - S_1$}{The maximum value of $x_1 + x_2$ is $3$ on the region $S_1 \cap S_2$}{The maximum value of $x_1 + x_2$ is $4$ on the region $S_1 \cup S_2$}
\ans{(C), (D) (MSQ)}

\textbf{Q.26 (MSQ)} Let $f : \mathbb{R}^2 - \{(0, 0)\} \to \mathbb{R}$ be a function defined by
\[
  f(x, y) = \left(\frac{x^2 - y^2}{x^2 + y^2}\right) + x \sin\!\left(\frac{1}{x^2 + y^2}\right).
\]
Consider the following three statements:\\
S1: $\lim_{x \to 0} \lim_{y \to 0} f(x, y)$ exists.\\
S2: $\lim_{y \to 0} \lim_{x \to 0} f(x, y)$ exists.\\
S3: $\lim_{(x, y) \to (0, 0)} f(x, y)$ exists.\\
Then, which of the following is/are correct?
\optslistA{S2 and S3 are TRUE and S1 is FALSE}{S1 and S2 are TRUE and S3 is FALSE}{S1 and S3 are TRUE and S2 is FALSE}{S1, S2 and S3 all are TRUE}
\ans{(B) (MSQ)}

\subsection*{Q.27 -- Q.35 Numerical Answer Type (NAT), carry ONE mark each (no negative marking).}

\textbf{Q.27} Let $M$ be a $7 \times 7$ matrix with entries in $\mathbb{R}$ and having the characteristic polynomial
\[
  c_M(x) = (x - 1)^\alpha (x - 2)^\beta (x - 3)^2, \quad \text{where } \alpha > \beta.
\]
Let $\operatorname{rank}(M - I_7) = \operatorname{rank}(M - 2I_7) = \operatorname{rank}(M - 3I_7) = 5$, where $I_7$ is the $7 \times 7$ identity matrix. If $m_M(x)$ is the minimal polynomial of $M$, then $m_M(5)$ is equal to \underline{\hspace{2cm}} (in integer).
\ans{$96$}

\textbf{Q.28} Let $y = P_n(x)$ be the unique polynomial of degree $n$ satisfying the Legendre differential equation
\[
  (1 - x^2)y'' - 2xy' + n(n + 1)y = 0 \quad \text{and} \quad y(1) = 1.
\]
Then, the value of $P_{11}'(1)$ is equal to \underline{\hspace{2cm}} (in integer).
\ans{$66$}

\textbf{Q.29} Let $\hat{\mathbf{a}}$ be a unit vector parallel to the tangent at the point $P(1, 1, \sqrt{2})$ to the curve of intersection of the surfaces $2x^2 + 3y^2 - z^2 = 3$ and $x^2 + y^2 = z^2$. Then, the absolute value of the directional derivative of
\[
  f(x, y, z) = x^2 + 2y^2 - 2\sqrt{11}\,z
\]
at $P$ in the direction of $\hat{\mathbf{a}}$ is \underline{\hspace{2cm}} (in integer).
\ans{$2$}

\textbf{Q.30} The volume of the region bounded by the cylinders $x^2 + y^2 = 4$ and $x^2 + z^2 = 4$ is \underline{\hspace{2cm}} (rounded off to TWO decimal places).
\ans{$42.50 \text{ to } 42.70$}

\textbf{Q.31} Let $W$ be the vector space (over $\mathbb{R}$) consisting of all bounded real-valued solutions of the differential equation
\[
  \frac{d^4 y}{dx^4} + 2\frac{d^2 y}{dx^2} + y = 0.
\]
Then, the dimension of $W$ is \underline{\hspace{2cm}} (in integer).
\ans{$2$}

\textbf{Q.32} Let $\mathbf{F} = (y - z)\hat{\mathbf{i}} + (z - x)\hat{\mathbf{j}} + (x - y)\hat{\mathbf{k}}$ be a vector field, and let $S$ be the surface $x^2 + y^2 + (z - 1)^2 = 9$, $1 \le z \le 4$. If $\hat{\mathbf{n}}$ denotes the unit outward normal vector to $S$, then the value of
\[
  \frac{1}{\pi} \left|\iint_S (\nabla \times \mathbf{F}) \cdot \hat{\mathbf{n}}\,dS\right|
\]
is equal to \underline{\hspace{2cm}} (in integer).
\ans{$18$}

\textbf{Q.33} Consider
\[
  I = \frac{1}{2\pi i} \oint_C \frac{\sin z}{1 - \cos(z^3)}\,dz,
\]
where $C = \{z \in \mathbb{C} : z = x + iy, \; |x| + |y| = 1, \; x, y \in \mathbb{R}\}$ is oriented positively as a simple closed curve. Then, the value of $120I$ is equal to \underline{\hspace{2cm}} (in integer).
\ans{$2$}

\textbf{Q.34} Let $\alpha, \beta, \gamma, \delta \in \mathbb{R}$ be such that the quadrature formula
\[
  \int_{-1}^{1} f(x)\,dx = \alpha f(-1) + \beta f(1) + \gamma f'(-1) + \delta f'(1)
\]
is exact for all polynomials of degree less than or equal to $3$. Then, $9(\alpha^2 + \beta^2 + \gamma^2 + \delta^2)$ is equal to \underline{\hspace{2cm}} (in integer).
\ans{$20$}

\textbf{Q.35} Let $y(x)$ be the solution of the initial value problem
\[
  \frac{dy}{dx} = \sin(\pi(x + y)), \quad y(0) = 0.
\]
Using Euler's method, with the step-size $h = 0.5$, the approximate value of $y(1.5) + 2y(1)$ is equal to \underline{\hspace{2cm}} (in integer).
\ans{$1$}

\subsection*{Q.36 -- Q.47 Multiple Choice Question (MCQ), carry TWO marks each (for each wrong answer: $-2/3$).}

\textbf{Q.36} Consider the linear system $A\mathbf{x} = \mathbf{b}$, where $A = [a_{ij}]$, $i, j = 1, 2, 3$, and $a_{ii} \ne 0$ for $i = 1, 2, 3$, is a matrix with entries in $\mathbb{R}$. For $D = \begin{bmatrix} a_{11} & 0 & 0 \\ 0 & a_{22} & 0 \\ 0 & 0 & a_{33} \end{bmatrix}$, let
\[
  D^{-1}A = \begin{bmatrix} 1 & 1 & -2 \\ 3 & 1 & 2 \\ 1 & 1 & 1 \end{bmatrix} \quad \text{and} \quad D^{-1}\mathbf{b} = \begin{bmatrix} 4 \\ 4 \\ 1 \end{bmatrix}.
\]
Consider the following two statements:\\
S1: The approximation of $\mathbf{x}$ after one iteration of the Jacobi scheme with initial vector $\mathbf{x}_0 = [1 \;\; 1 \;\; 1]^T$ is $\mathbf{x}_1 = [5 \;\; -1 \;\; -1]^T$.\\
S2: There exists an initial vector $\mathbf{x}_0$ for which Jacobi iterative scheme diverges.\\
Then, which one of the following is correct?
\optslistA{S1 is TRUE and S2 is FALSE}{S2 is TRUE and S1 is FALSE}{both S1 and S2 are TRUE}{neither S1 nor S2 is TRUE}
\ans{(C)}

\textbf{Q.37} Let $y(x)$ be the solution of the differential equation
\[
  x^2 y'' + 7x y' + 9y = x^{-3} \log_e x, \quad x > 0,
\]
satisfying $y(1) = 0$ and $y'(1) = 0$. Then, the value of $y(e)$ is equal to
\optsfourA{$\dfrac{1}{3} e^{-3}$}{$\dfrac{1}{6} e^{-3}$}{$\dfrac{2}{3} e^{-3}$}{$\dfrac{1}{2} e^{-3}$}
\ans{(B)}

\textbf{Q.38} Let $y_1(x)$ and $y_2(x)$ be the two linearly independent solutions of the differential equation
\[
  (1 + x^2) y'' - x y' + (\cos^2 x) y = 0,
\]
satisfying the initial conditions $y_1(0) = 3$, $y_1'(0) = -1$ and $y_2(0) = -5$, $y_2'(0) = 2$. Define $W(x) = \begin{vmatrix} y_1(x) & y_2(x) \\ y_1'(x) & y_2'(x) \end{vmatrix}$. Then, the value of $W\!\left(\frac{1}{2}\right)$ is
\optsfourA{$\dfrac{\sqrt{5}}{4}$}{$\dfrac{\sqrt{5}}{2}$}{$\dfrac{2}{\sqrt{5}}$}{$\dfrac{4}{\sqrt{5}}$}
\ans{(B)}

\textbf{Q.39} Let $C$ be the curve of intersection of the surfaces $z^2 = x^2 + y^2$ and $4x + z = 7$. If $P$ is a point on $C$ at a minimum distance from the $xy$-plane, then the distance of $P$ from the origin is
\optsfourA{$\dfrac{7}{5}$}{$\dfrac{7\sqrt{2}}{5}$}{$\dfrac{14}{5}$}{$\dfrac{14\sqrt{2}}{5}$}
\ans{(B)}

\textbf{Q.40} Let $u(x, t)$ be the solution of the initial-value problem
\[
  \frac{\partial^2 u}{\partial t^2} - 9\frac{\partial^2 u}{\partial x^2} = 0, \quad x \in \mathbb{R}, \; t > 0, \qquad u(x, 0) = e^x, \quad \frac{\partial u}{\partial t}(x, 0) = \sin x.
\]
Then, the value of $u\!\left(\frac{\pi}{2}, \frac{\pi}{6}\right)$ is
\optsfourA{$\dfrac{1}{2}\left(e^\pi - \dfrac{1}{3}\right)$}{$\dfrac{1}{2}\left(e^\pi + \dfrac{1}{3}\right)$}{$\dfrac{1}{2}\left(e^\pi + \dfrac{5}{3}\right)$}{$\dfrac{1}{2}\left(e^\pi - \dfrac{5}{3}\right)$}
\ans{(C)}

\textbf{Q.41} Let $T$ be the M\"obius transformation that maps the points $0, \frac{1}{2}$ and $1$ conformally onto the points $-3, \infty$ and $2$, respectively, in the extended complex plane. If $T$ maps the circle centred at $1$ with radius $k$ onto a straight line given by the equation $\alpha x + \beta y + \gamma = 0$, then the value of $\dfrac{2k(\alpha + \beta) + \gamma}{\alpha + \beta - 2k\gamma}$ is equal to
\optsfourA{$\dfrac{1}{7}$}{$\dfrac{2}{7}$}{$\dfrac{1}{3}$}{$\dfrac{2}{3}$}
\ans{(A)}

\textbf{Q.42} Let $U = \{z \in \mathbb{C} : \operatorname{Im}(z) > 0\}$ and $D = \{z \in \mathbb{C} : |z| < 1\}$, where $\operatorname{Im}(z)$ denotes the imaginary part of $z$. Let $S$ be the set of all bijective analytic functions $f : U \to D$ such that $f(i) = 0$. Then, the value of $\sup_{f \in S} |f(4i)|$ is
\optsfourA{$0$}{$\dfrac{1}{4}$}{$\dfrac{1}{2}$}{$\dfrac{3}{5}$}
\ans{(D)}

\textbf{Q.43} Let $\Omega$ be a non-empty open connected subset of $\mathbb{C}$ and $f : \Omega \to \mathbb{C}$ be a non-constant function. Let the functions $f^2 : \Omega \to \mathbb{C}$ and $f^3 : \Omega \to \mathbb{C}$ be defined by $f^2(z) = (f(z))^2$ and $f^3(z) = (f(z))^3$, $z \in \Omega$. Consider the following two statements:\\
S1: If $f$ is continuous in $\Omega$ and $f^2$ is analytic in $\Omega$, then $f$ is analytic in $\Omega$.\\
S2: If $f^2$ and $f^3$ are analytic in $\Omega$, then $f$ is analytic in $\Omega$.\\
Then, which one of the following is correct?
\optslistA{S1 is TRUE and S2 is FALSE}{S2 is TRUE and S1 is FALSE}{both S1 and S2 are TRUE}{neither S1 nor S2 is TRUE}
\ans{(C)}

\textbf{Q.44} In the following, all subsets of Euclidean spaces are considered with the respective subspace topologies. Define an equivalence relation $\sim$ on the sphere
\[
  S = \{(x_1, x_2, x_3) \in \mathbb{R}^3 : x_1^2 + x_2^2 + x_3^2 = 1\}
\]
by $(x_1, x_2, x_3) \sim (y_1, y_2, y_3)$ if $x_3 = y_3$, for $(x_1, x_2, x_3), (y_1, y_2, y_3) \in S$. Let $[x_1, x_2, x_3]$ denote the equivalence class of $(x_1, x_2, x_3)$, and let $X$ denote the set of all such equivalence classes. Let $\mathcal{L} : S \to X$ be given by $\mathcal{L}((x_1, x_2, x_3)) = [x_1, x_2, x_3]$. If $X$ is provided with the quotient topology induced by the map $\mathcal{L}$, then which one of the following is TRUE?
\optslistA{$X$ is homeomorphic to $\{x \in \mathbb{R} : -1 \le x \le 1\}$}{$X$ is homeomorphic to $\{(x_1, x_2) \in \mathbb{R}^2 : x_1^2 + x_2^2 = 1\}$}{$X$ is homeomorphic to $\{(x_1, x_2) \in \mathbb{R}^2 : x_1^2 + x_2^2 \le 1\}$}{$X$ is homeomorphic to $\{(x_1, x_2, x_3) \in \mathbb{R}^3 : x_1^2 + x_2^2 = 1 \text{ and } -1 \le x_3 \le 1\}$}
\ans{(A)}

\textbf{Q.45} Consider the following two spaces:\\
$X = (C[-1, 1], \|\cdot\|_\infty)$, the space of all real-valued continuous functions defined on $[-1, 1]$ equipped with the norm $\|f\|_\infty = \sup_{t \in [-1, 1]} |f(t)|$.\\
$Y = (C[-1, 1], \|\cdot\|_2)$, the space of all real-valued continuous functions defined on $[-1, 1]$ equipped with the norm $\|f\|_2 = (\int_{-1}^{1} |f(t)|^2\,dt)^{1/2}$.\\
Let $W$ be the linear span over $\mathbb{R}$ of all the Legendre polynomials. Then, which one of the following is correct?
\optslistA{$W$ is dense in $X$ but not in $Y$}{$W$ is dense in $Y$ but not in $X$}{$W$ is dense in both $X$ and $Y$}{$W$ is dense neither in $X$ nor in $Y$}
\ans{(C)}

\textbf{Q.46} Consider the metric spaces $X = (\mathbb{R}, d_1)$ and $Y = ([0, 1], d_2)$ with the metrics defined by $d_1(x, y) = |x - y|$, $x, y \in \mathbb{R}$ and $d_2(x, y) = |x - y|$, $x, y \in [0, 1]$, respectively. Then, which one of the following is TRUE?
\optslistA{$[0, \frac{1}{4})$ is open in $X$ but not in $Y$}{$[0, \frac{1}{4})$ is open in $Y$ but not in $X$}{$[0, \frac{1}{4})$ is open in both $X$ and $Y$}{$[0, \frac{1}{4})$ is open neither in $X$ nor in $Y$}
\ans{(B)}

\textbf{Q.47} Let $K$ be an algebraically closed field containing a finite field $F$. Let $L$ be the subfield of $K$ consisting of elements of $K$ that are algebraic over $F$. Consider the following statements:\\
S1: $L$ is algebraically closed.\\
S2: $L$ is infinite.\\
Then, which one of the following is correct?
\optslistA{S1 is TRUE and S2 is FALSE}{S2 is TRUE and S1 is FALSE}{both S1 and S2 are TRUE}{neither S1 nor S2 is TRUE}
\ans{(C)}

\subsection*{Q.48 -- Q.60 Multiple Select Question (MSQ), carry TWO marks each (no negative marking).}

\textbf{Q.48 (MSQ)} Let $M_2(\mathbb{R})$ be the vector space (over $\mathbb{R}$) of all $2 \times 2$ matrices with entries in $\mathbb{R}$. Consider the linear transformation $T : M_2(\mathbb{R}) \to M_2(\mathbb{R})$ defined by $T(X) = AXB$, where $A = \begin{bmatrix} 1 & -2 \\ 1 & 4 \end{bmatrix}$ and $B = \begin{bmatrix} 6 & 5 \\ -2 & -1 \end{bmatrix}$. If $P$ is the matrix representation of $T$ with respect to the standard basis of $M_2(\mathbb{R})$, then which of the following is/are TRUE?
\optslistA{$P$ is an invertible matrix}{The trace of $P$ is $25$}{The rank of $(P^2 - 4I_4)$ is $4$, where $I_4$ is the $4 \times 4$ identity matrix}{The nullity of $(P - 2I_4)$ is $0$, where $I_4$ is the $4 \times 4$ identity matrix}
\ans{(A), (B) (MSQ)}

\textbf{Q.49 (MSQ)} Consider the linear programming problem (LPP):
\begin{align*}
  \text{Maximize} \quad & Z = 3x_1 + 5x_2 \\
  \text{Subject to} \quad & x_1 + x_3 = 4, \\
  & 2x_2 + x_4 = 12, \\
  & 3x_1 + 2x_2 + x_5 = 18, \\
  & x_1, x_2, x_3, x_4, x_5 \ge 0.
\end{align*}
Given that $x_B = (x_3, x_2, x_1)^T$ forms the optimal basis of the LPP with basis matrix $B$ and respective $B^{-1} = \begin{bmatrix} \alpha & \beta & -\beta \\ 0 & \gamma & 0 \\ 0 & -\beta & \beta \end{bmatrix}$. If $(p, q, r)$ is the optimal solution of the dual of the LPP, then which of the following is/are TRUE?
\optsfourA{$\alpha + 3\beta + 2\gamma = 3$}{$\alpha - 3\beta + 4\gamma = 1$}{$p + q + r = \dfrac{5}{2}$}{$p^2 + q^2 + r^2 = \dfrac{17}{4}$}
\ans{(A), (C) (MSQ)}

\textbf{Q.50 (MSQ)} Let $0 < \alpha < 1$. Define
\[
  C^\alpha[0, 1] = \left\{f : [0, 1] \to \mathbb{R} : \sup_{\substack{s \ne t \\ s, t \in [0, 1]}} \frac{|f(t) - f(s)|}{|t - s|^\alpha} < \infty\right\}.
\]
It is given that $C^\alpha[0, 1]$ is a Banach space with respect to the norm $\|\cdot\|_\alpha$ given by
\[
  \|f\|_\alpha = |f(0)| + \sup_{\substack{s \ne t \\ s, t \in [0, 1]}} \frac{|f(t) - f(s)|}{|t - s|^\alpha}.
\]
Let $C[0, 1]$ be the space of all real-valued continuous functions on $[0, 1]$ with the norm $\|f\|_\infty = \sup_{0 \le t \le 1} |f(t)|$. If $T : C^\alpha[0, 1] \to C[0, 1]$ is the map $Tf = f$, $f \in C^\alpha[0, 1]$, then which of the following is/are TRUE?
\optslistA{$T$ is a compact linear map}{Image of $T$ is closed in $C[0, 1]$}{Image of $T$ is dense in $C[0, 1]$}{$T$ is not a bounded linear map}
\ans{(A), (C) (MSQ)}

\textbf{Q.51 (MSQ)} Let $u(x, t)$ be the solution of the initial value problem
\[
  \frac{\partial u}{\partial t} + 3\frac{\partial u}{\partial x} = u, \quad x \in \mathbb{R}, \; t > 0, \qquad u(x, 0) = \cos x,
\]
and let $v(x, t)$ be the solution of the initial value problem
\[
  \frac{\partial v}{\partial t} + 3\frac{\partial v}{\partial x} = v^2, \quad x \in \mathbb{R}, \; t > 0, \qquad v(x, 0) = \cos x.
\]
Then, which of the following is/are TRUE?
\optslistA{$|u(x, t)| \le e^t$ for all $x \in \mathbb{R}$ and for all $t > 0$}{$v(x, 1)$ is not defined for certain values of $x \in \mathbb{R}$}{$v(x, 1)$ is not defined for any $x \in \mathbb{R}$}{$u(2\pi, \pi) = -e^\pi$}
\ans{(A), (B), (D) (MSQ)}

\textbf{Q.52 (MSQ)} Let $u(x, t)$ be the solution of the initial-boundary value problem
\[
  \frac{\partial u}{\partial t} = 2\frac{\partial^2 u}{\partial x^2}, \quad 0 < x < 1, \; t > 0, \qquad u(0, t) = u(1, t) = 0, \quad u(x, 0) = 2x(1 - x).
\]
Then, which of the following is/are TRUE?
\optslistA{$0 \le u(x, t) \le \frac{1}{4}$ for all $t \ge 0$ and $x \in [0, 1]$}{$u(x, t) = u(1 - x, t)$ for all $t \ge 0$ and $x \in [0, 1]$}{$\int_0^1 (u(x, t))^2\,dx$ is a decreasing function of $t$}{$\int_0^1 (u(x, t))^2\,dx$ is not a decreasing function of $t$}
\ans{(B), (C) (MSQ)}

\textbf{Q.53 (MSQ)} Consider the function $f : \mathbb{R}^2 \to \mathbb{R}^2$ given by
\[
  f(x, y) = (e^{2\pi x} \cos 2\pi y, \; e^{2\pi x} \sin 2\pi y).
\]
Then, which of the following is/are TRUE?
\optslistA{If $G$ is open in $\mathbb{R}^2$, then $f(G)$ is open in $\mathbb{R}^2$}{If $G$ is closed in $\mathbb{R}^2$, then $f(G)$ is closed in $\mathbb{R}^2$}{If $G$ is dense in $\mathbb{R}^2$, then $f(G)$ is dense in $\mathbb{R}^2$}{$f$ is surjective}
\ans{(A), (C) (MSQ)}

\textbf{Q.54 (MSQ)} Let $\{x_k\}_{k=1}^{\infty}$ be an orthonormal set of vectors in a real Hilbert space $X$ with inner product $\langle \cdot, \cdot \rangle$. Let $n \in \mathbb{N}$, and let $Y$ be the linear span of $\{x_k\}_{k=1}^{n}$ over $\mathbb{R}$. For $x \in X$, let $S_n(x) = \sum_{k=1}^{n} \langle x, x_k \rangle x_k$. Then, which of the following is/are TRUE?
\optslistA{$S_n(x)$ is the orthogonal projection of $x$ onto $Y$}{$S_n(x)$ is the orthogonal projection of $x$ onto $Y^\perp$}{$(x - S_n(x))$ is orthogonal to $S_n(x)$ for all $x$ in $X$}{$\sum_{k=1}^{n} |\langle x, x_k \rangle|^2 = \|x\|^2$ for all $x$ in $X$}
\ans{(A), (C) (MSQ)}

\textbf{Q.55 (MSQ)} Consider the sequence $\{f_n\}$ of continuous functions on $[0, 1]$ defined by
\[
  f_1(x) = \frac{x}{2}, \quad f_{n+1}(x) = f_n(x) - \frac{1}{2}\bigl((f_n(x))^2 - x\bigr), \quad n = 1, 2, 3, \ldots
\]
Then, which of the following is/are TRUE?
\optslistA{The sequence $\{f_n\}$ converges pointwise but not uniformly on $[0, 1]$}{The sequence $\{f_n\}$ converges uniformly on $[0, 1]$}{$\sqrt{x} - f_n(x) > \dfrac{2\sqrt{x}}{2 + n\sqrt{x}}$ for all $x \in [0, 1]$ and $n = 1, 2, 3, \ldots$}{$0 \le f_n(x) \le \sqrt{x}$ for all $x \in [0, 1]$ and $n = 1, 2, 3, \ldots$}
\ans{(B), (D) (MSQ)}

\textbf{Q.56 (MSQ)} For $x \in (0, \pi)$, let $u_n(x) = \dfrac{\sin nx}{\sqrt{n}}$, $n = 1, 2, 3, \ldots$. Then, which of the following is/are TRUE?
\optslistA{$\sum_{n=1}^{\infty} u_n(x)$ converges uniformly on $(0, \pi)$}{$\sum_{n=1}^{\infty} u_n'(x)$ converges uniformly on $(0, \pi)$}{$\sum_{n=1}^{\infty} u_n(x)$ converges pointwise but not uniformly on $(0, \pi)$}{$\sum_{n=1}^{\infty} u_n(x)$ converges uniformly on every compact subset of $(0, \pi)$}
\ans{(C), (D) (MSQ)}

\textbf{Q.57 (MSQ)} Let $\mathbb{R}^1$ and $\mathbb{R}^2$ be provided with the respective Euclidean topologies, and let $S^1 = \{(x_1, x_2) \in \mathbb{R}^2 : x_1^2 + x_2^2 = 1\}$ be assigned the subspace topology induced from $\mathbb{R}^2$. If $f : S^1 \to \mathbb{R}^1$ is a non-constant continuous function, then which of the following is/are TRUE?
\optslistA{$f$ maps closed sets to closed sets}{$f$ is injective}{$f$ is surjective}{There exists $\lambda \in \mathbb{R}$ such that $f(\cos\lambda, \sin\lambda) = f(-\cos\lambda, -\sin\lambda)$}
\ans{(A), (D) (MSQ)}

\textbf{Q.58 (MSQ)} Let $X$ be an uncountable set. Let the topology on $X$ be defined by declaring a subset $U \subseteq X$ to be open if $X - U$ is either empty or finite or countable, and the empty set to be open. Then, which of the following is/are TRUE?
\optslistA{Every compact subset of $X$ is closed}{Every closed subset of $X$ is compact}{$X$ is $T_1$ (singleton subsets are closed) but not $T_2$ (Hausdorff)}{$X$ is $T_2$ (Hausdorff)}
\ans{(A), (C) (MSQ)}

\textbf{Q.59 (MSQ)} All rings considered below are assumed to be associative and commutative with $1 \ne 0$. Further, all ring homomorphisms map $1$ to $1$. Consider the following statements about such a ring $R$:\\
P1: $R$ is isomorphic to the product of two rings $R_1$ and $R_2$.\\
P2: $\exists \, r_1, r_2 \in R$ such that $r_1^2 = r_1 \ne 0 \ne r_2 = r_2^2$, $r_1 r_2 = 0$ and $r_1 + r_2 = 1$.\\
P3: $\exists$ ideals $I_1, I_2 \subseteq R$ with $R \ne I_1 \ne (0) \ne I_2 \ne R$ such that $R = I_1 + I_2$ and $I_1 \cap I_2 = (0)$.\\
P4: $\exists \, a, b \in R$ with $a \ne 0 \ne b$ such that $ab = 0$.\\
Then, which of the following is/are TRUE?
\optsfourA{$P1 \implies P2$}{$P2 \implies P3$}{$P3 \implies P4$}{$P4 \implies P1$}
\ans{(A), (B), (C) (MSQ)}

\textbf{Q.60 (MSQ)} Let $E \subset F$ and $F \subset K$ be field extensions which are not algebraic. Let $\alpha \in K$ be algebraic over $F$ and $\alpha \notin F$. Let $L$ be the subfield of $K$ generated over $E$ by the coefficients of the monic polynomial of minimal degree over $F$ which has $\alpha$ as a zero. Then, which of the following is/are TRUE?
\optslistA{$F(\alpha) \supset L(\alpha)$ is a finite extension if and only if $F \supset L$ is a finite extension}{The dimension of $L(\alpha)$ over $L$ is greater than the dimension of $F(\alpha)$ over $F$}{The dimension of $L(\alpha)$ over $L$ is smaller than the dimension of $F(\alpha)$ over $F$}{$F(\alpha) \supset L(\alpha)$ is an algebraic extension if and only if $F \supset L$ is an algebraic extension}
\ans{(A), (D) (MSQ)}

\subsection*{Q.61 -- Q.65 Numerical Answer Type (NAT), carry TWO marks each (no negative marking).}

\textbf{Q.61} Consider the inner product space of all real-valued continuous functions defined on $[-1, 1]$ with the inner product $\langle f, g \rangle = \int_{-1}^{1} f(x) g(x)\,dx$. If $p(x) = \alpha + \beta x^2 - 30x^4$, $\alpha, \beta \in \mathbb{R}$, is orthogonal to all the polynomials having degree less than or equal to $3$, with respect to this inner product, then $\alpha + 5\beta$ is equal to \underline{\hspace{2cm}} (in integer).
\ans{$126$}

\textbf{Q.62} For $X = (x_1, x_2, x_3)^T \in \mathbb{R}^3$, consider the quadratic form:
\[
  Q(X) = 2x_1^2 + 2x_2^2 + 3x_3^2 + 4x_1 x_2 + 2x_1 x_3 + 2x_2 x_3.
\]
Let $M$ be the symmetric matrix associated with the quadratic form $Q(X)$ with respect to the standard basis of $\mathbb{R}^3$. Let $Y = (y_1, y_2, y_3)^T \in \mathbb{R}^3$ be a non-zero vector, and let
\[
  a_n = \frac{Y^T (M + I_3)^{n+1} Y}{Y^T (M + I_3)^n Y}, \quad n = 1, 2, 3, \ldots,
\]
where $I_3$ is the $3 \times 3$ identity matrix. Then, the value of $\lim_{n \to \infty} a_n$ is \underline{\hspace{2cm}} (in integer).
\ans{$6$}

\textbf{Q.63} Let $\alpha, \beta$ be distinct non-zero real numbers, and let $Q(z)$ be a polynomial of degree less than $5$. If the function
\[
  f(z) = \frac{\alpha^6 \sin \beta z - \beta^6(e^{2\alpha z} - Q(z))}{z^6}
\]
satisfies Morera's theorem in $\mathbb{C} \setminus \{0\}$, then the value of $\dfrac{\alpha}{4\beta}$ is equal to \underline{\hspace{2cm}} (in integer).
\ans{$8$}

\textbf{Q.64} Let $G$ be a group with identity element $e$, and let $g, h \in G$ be such that the following hold:
\begin{itemize}
  \setlength{\itemsep}{1pt}\setlength{\parskip}{0pt}
  \item[(i)] $g \ne e$, $g^2 = e$,
  \item[(ii)] $h \ne e$, $h^2 \ne e$, and $g h g^{-1} = h^2$.
\end{itemize}
Then, the least positive integer $n$ for which $h^n = e$ is \underline{\hspace{2cm}} (in integer).
\ans{$3$}

\textbf{Q.65} Let $(\mathbb{R}^2, d_1)$ and $(\mathbb{R}^2, d_2)$ be two metric spaces with
\[
  d_1((x_1, x_2), (y_1, y_2)) = |x_1 - y_1| + |x_2 - y_2| \quad \text{and} \quad d_2((x_1, x_2), (y_1, y_2)) = \frac{d_1((x_1, x_2), (y_1, y_2))}{1 + d_1((x_1, x_2), (y_1, y_2))}.
\]
If the open ball centred at $(0, 0)$ with radius $\frac{1}{7}$ in $(\mathbb{R}^2, d_1)$ is equal to the open ball centred at $(0, 0)$ with radius $\frac{1}{\alpha}$ in $(\mathbb{R}^2, d_2)$, then the value of $\alpha$ is \underline{\hspace{2cm}} (in integer).
\ans{$8$}

% ============================================================================
% GATE 2026 -- GENERAL APTITUDE (GA) & MATHEMATICS (MA)
% ============================================================================

\newpage
\begin{center}
  {\LARGE\bfseries GATE 2026}\\[4pt]
  {\large\bfseries General Aptitude (GA) \& Mathematics (MA)}\\[2pt]
  {\small\itshape Official Final Answer Key attached to every question}
\end{center}
\medskip

\section*{General Aptitude (GA)}

\subsection*{Q.1 -- Q.5 Multiple Choice Question (MCQ), carry ONE mark each (for each wrong answer: $-1/3$).}

\textbf{Q.1} Suresh said, ``I did it yesterday.''\\
Which one of the following options is the correct form of this sentence in indirect speech?
\optslistA{Suresh said that I did it yesterday.}{Suresh says I did it yesterday.}{Suresh says that he did it the day before.}{Suresh said that he had done it the day before.}
\ans{(D)}

\textbf{Q.2} To continue the sequence of tiles shown, the tile indicated by the question mark should be\\[4pt]
\emph{[Diagram note: Six square domino-style tiles show dot counts following the Fibonacci sequence $0, 1, 1, 2, 3, 5$, followed by a blank tile marked ``?''. The four options show tiles with $4$, $6$, $8$, and $9$ dots, respectively.]}
\optsfourA{Tile with $4$ dots}{Tile with $6$ dots}{Tile with $8$ dots}{Tile with $9$ dots}
\ans{(C)}

\textbf{Q.3} Consider an art gallery whose walkways are shown as lines in the diagram. A black dot represents a junction of two walkways. A guard may be placed at a junction to watch over the walkways that join at that junction. The minimum number of guards needed to watch all the walkways is \underline{\hspace{2cm}}.\\
\emph{[Diagram note: A pentagon (5-cycle $C_5$ shaped like a house outline) with $5$ walkways (edges) and $5$ black-dot junctions (vertices).]}
\optsfourA{$2$}{$3$}{$4$}{$5$}
\ans{(B)}

\textbf{Q.4} The 2nd of June is a Thursday in a certain year. Which day of the week is the 3rd of July in that year?
\optsfourA{Thursday}{Friday}{Saturday}{Sunday}
\ans{(D)}

\textbf{Q.5} A coin with heads facing up is shown as \textsf{H} and a coin with tails facing up is shown as \textsf{T}. Six coins are placed in the Starting Arrangement, as shown below:
\[
  \text{Starting Arrangement: } \textsf{H}\;\textsf{H}\;\textsf{H}\;\textsf{T}\;\textsf{T}\;\textsf{T} \qquad \longrightarrow \qquad \text{Final Arrangement: } \textsf{T}\;\textsf{T}\;\textsf{T}\;\textsf{H}\;\textsf{H}\;\textsf{H}
\]
A ``step'' is defined as interchanging a pair of adjacent coins without flipping them. The minimum number of steps needed to go from the Starting Arrangement to the Final Arrangement is \underline{\hspace{2cm}}.
\optsfourA{$3$}{$6$}{$9$}{$12$}
\ans{(C)}

\subsection*{Q.6 -- Q.10 Multiple Choice Question (MCQ), carry TWO marks each (for each wrong answer: $-2/3$).}

\textbf{Q.6} Exacerbate : Mitigate :: \underline{\hspace{2cm}}\\
Choose the option with the correct pair of words to fill the blank.
\optsfourA{Aggravate : Alleviate}{Alleviate : Precipitate}{Aggravate : Precipitate}{Emancipate : Exonerate}
\ans{(A)}

\textbf{Q.7} A paper shown in Panel I is folded along the dashed lines (- - -) to construct a cube. The shaded regions shown in Panel I appear on the outer surface of the cube. Referring to cubes shown in Panel II, which one of the options is correct?\\[4pt]
\emph{[Diagram note: Panel I shows a cross-shaped 6-face cube net (4 faces in a vertical column plus 1 face to the left and 1 face to the right of the second face from the top): top face has a grey bottom-right right-triangle; second face has a thin vertical grey stripe down the middle; left face has a grey L-shape along its left and bottom edges; right face has a grey small square in its top-right corner; third face has a thin horizontal grey stripe across the middle; bottom face is half-shaded grey along its bottom half. Panel II shows two 3D cube views: (i) front face with an inverted-T/L grey block, top face with a horizontal stripe, right face with a half-triangle; (ii) front face with a grey bottom-right right-triangle, top face with a top-edge grey band, right face with a bottom-right square/corner shading.]}
\optslistA{Only (i) can correspond to the unfolded cube in Panel I.}{Only (ii) can correspond to the unfolded cube in Panel I.}{Both (i) and (ii) can correspond to the unfolded cube in Panel I.}{Neither (i) nor (ii) can correspond to the unfolded cube in Panel I.}
\ans{(B)}

\textbf{Q.8} In a population, patients who have high cholesterol also have high blood-pressure (BP). Some patients with high BP also have diabetes. There are no patients who have both high cholesterol and diabetes. Furthermore,
\begin{enumerate}
  \setlength{\itemsep}{1pt}\setlength{\parskip}{0pt}
  \item the total number of patients with at least one of these conditions is $75$,
  \item the number of patients with high cholesterol is $10$,
  \item the number of patients with high BP is $45$, and
  \item the number of patients with only high BP and no other conditions is $20$.
\end{enumerate}
Then the number of patients who have both diabetes and high BP is \underline{\hspace{2cm}}.
\optsfourA{$0$}{$15$}{$20$}{$10$}
\ans{(B)}

\textbf{Q.9} Four people P, Q, R, and S, of different ages, make the following observations:\\
P -- I am younger than S.\\
Q -- I am neither the youngest nor the oldest.\\
R -- P is older than me.\\
Based on these observations, the youngest person is \underline{\hspace{2cm}}.
\optsfourA{P}{Q}{R}{S}
\ans{(C)}

\textbf{Q.10} Circles $C_1, C_2$, and $C_3$, with centers $O_1, O_2$, and $O_3$, and radii $r_1, r_2$, and $r_3$, respectively, touch each other externally. Given $r_1 = 2\text{ cm}$, $r_2 = 1\text{ cm}$ and the angle $\angle O_1 O_3 O_2$ is $90^\circ$, $r_3 = \underline{\hspace{1.5cm}}\text{ cm}$.
\optsfourA{$\dfrac{1}{2}(-3 + \sqrt{17})$}{$\dfrac{1}{2}(3 + \sqrt{17})$}{$\dfrac{1}{2}(-2 + \sqrt{17})$}{$\dfrac{1}{2}(-3 + 2\sqrt{17})$}
\ans{(A)}

\section*{Mathematics (MA)}

\subsection*{Q.11 -- Q.20 Multiple Choice Question (MCQ), carry ONE mark each (for each wrong answer: $-1/3$).}

\textbf{Q.11} Let $M_3(\mathbb{R})$ be the vector space of all $3 \times 3$ real matrices over $\mathbb{R}$ under usual matrix addition and scalar multiplication. Let $T_3(\mathbb{R})$ be the set of all $3 \times 3$ real upper triangular matrices. Which one of the following is TRUE?
\optslistA{The quotient space $M_3(\mathbb{R})/T_3(\mathbb{R})$ is isomorphic to the vector space of all $3 \times 3$ real skew symmetric matrices over $\mathbb{R}$ under usual matrix addition and scalar multiplication.}{The quotient space $M_3(\mathbb{R})/T_3(\mathbb{R})$ is isomorphic to the vector space of all $3 \times 3$ real symmetric matrices over $\mathbb{R}$ under usual matrix addition and scalar multiplication.}{The quotient space $M_3(\mathbb{R})/T_3(\mathbb{R})$ is isomorphic to the vector space of all $3 \times 3$ real lower triangular matrices over $\mathbb{R}$ under usual matrix addition and scalar multiplication.}{The quotient space $M_3(\mathbb{R})/T_3(\mathbb{R})$ is isomorphic to the vector space of all $3 \times 3$ real matrices with trace zero over $\mathbb{R}$ under usual matrix addition and scalar multiplication.}
\ans{(A)}

\textbf{Q.12} Let $I$ be the integral defined as follows:
\[
  I = \int_0^1 \int_0^{\sqrt{y}} dx\,dy + \int_1^2 \int_{\sqrt{y-1}}^{1} dx\,dy.
\]
If the order of the integration is changed, then which one of the following is the correct expression for $I$?
\optslistA{$\displaystyle\int_0^{1/2} \int_{-x}^{x} dy\,dx + \int_{1/2}^{1} \int_{-x^2}^{-x^2+1} dy\,dx$}{$\displaystyle\int_0^1 \int_{x^2}^{x^2+1} dy\,dx$}{$\displaystyle\int_0^{1/2} \int_{x^2}^{x^2+1} dy\,dx + \int_{1/2}^{1} \int_{\sqrt{x}}^{\sqrt{x+1}} dy\,dx$}{$\displaystyle\int_0^1 \int_{x^2}^{x^2-1} dy\,dx$}
\ans{(B)}

\textbf{Q.13} Let $C$ denote the Cantor set and $f : [0, 1] \to \mathbb{R}$ be defined as follows:
\[
  f(x) = \begin{cases} x^{2026}, & \text{for } x \in C, \\ \cos(\pi x), & \text{for } x \in \left[0, \frac{1}{2}\right] \setminus C, \\ \sin(\pi x), & \text{for } x \in \left[\frac{1}{2}, 1\right] \setminus C. \end{cases}
\]
The value of the Lebesgue integral of $f(x)$ over the interval $[0, 1]$ is equal to
\optsfourA{$\dfrac{2}{\pi}$}{$\dfrac{1}{\pi}$}{$\dfrac{3}{\pi}$}{$0$}
\ans{(A)}

\textbf{Q.14} Let $G$ be a group of order $595$. Which one of the following is TRUE?
\optslistA{$G$ cannot have a proper normal subgroup.}{$G$ must have a proper normal subgroup.}{$G$ cannot have an element of order $17$.}{The number of Sylow $5$-subgroups is $17$.}
\ans{(B)}

\textbf{Q.15} Let $u(x, t)$ be the solution of the initial value problem for the heat equation on the real line:
\[
  \frac{\partial u}{\partial t} = k \frac{\partial^2 u}{\partial x^2}, \quad -\infty < x < \infty, \; t > 0, \; k \in \mathbb{R},
\]
with the initial condition $u(x, 0) = e^{-a|x|}$, $a > 0$. If $\hat{u}(w, t) = \int_{-\infty}^{\infty} u(x, t) e^{iwx}\,dx$ is the Fourier transform of $u(x, t)$ with respect to $x$, then $\hat{u}(w, t)$ is equal to
\optsfourA{$\dfrac{2a}{a^2 + w^2} e^{-kw^2 t}$}{$\dfrac{a}{a^2 + w^2} e^{-kw^2 t}$}{$\dfrac{2a}{a^2 + w^2} e^{-ka^2 t}$}{$\dfrac{1}{\sqrt{4\pi kt}} e^{-w^2/(4kt)}$}
\ans{(A)}

\textbf{Q.16} Let $f : \mathbb{R}^2 \to \mathbb{R}^2$ be defined as $f(x, y) = (e^x \cos y, \, e^x \sin y)$. Which one of the following is TRUE?
\optsfourA{$f$ is one-to-one.}{The Jacobian of $f$ is negative.}{$f$ is locally invertible.}{$f$ is invertible on $\mathbb{R}^2$.}
\ans{(C)}

\textbf{Q.17} Let $f : \mathbb{C} \to \mathbb{C}$ be defined by $f(z) = |z|^2 - 5\overline{z} + 2$. Which one of the following is TRUE?
\optsfourA{$f$ is differentiable at $z = 5i$.}{$f$ is differentiable at $z = 5$.}{$f$ is differentiable at $z = -5i$.}{$f$ is differentiable at $z = -5$.}
\ans{(B)}

\textbf{Q.18} Let $D = \{z \in \mathbb{C} \mid |z| < 1\}$ denote the unit disc in the complex plane $\mathbb{C}$. Let $f : D \to D$ be an analytic function which satisfies $f(0) = 0$. Then which one of the following is a possible value of $f'(0)$?
\optsfourA{$\dfrac{5}{2}i$}{$\dfrac{i}{10}$}{$\dfrac{3}{2}$}{$-\dfrac{5}{2}i$}
\ans{(B)}

\textbf{Q.19} Let $X$ be the space of all continuously differentiable real valued functions on $[0, 1]$. Define the following norms on $X$:
\[
  p_1(x) = \sup\{|x(t)| : t \in [0, 1]\}, \quad p_2(x) = \sup\left\{\left|\frac{d}{dt} x(t)\right| : t \in [0, 1]\right\}, \quad \text{and} \quad p_3(x) = p_1(x) + p_2(x).
\]
Which one of the following is TRUE?
\optsfourA{$(X, p_1)$ is a Banach space.}{$(X, p_2)$ is a Banach space.}{$(X, p_3)$ is not a Banach space.}{$(X, p_3)$ is a Banach space.}
\ans{(D)}

\textbf{Q.20} Let $X$ be a set with at least two elements. Let $d_1, d_2$ and $d_3$ be metrics on $X$. Which one of the following is NOT a metric on $X$?
\optslistA{$d(x, y) := \min(d_1(x, y), 3)$ for all $x, y \in X$.}{$d(x, y) := \max(d_2(x, y), 3)$ for all $x, y \in X$.}{$d(x, y) := \dfrac{10\,d_3(x, y)}{1 + d_3(x, y)}$ for all $x, y \in X$.}{$d(x, y) := \dfrac{1}{3}(d_1(x, y) + d_2(x, y) + d_3(x, y))$ for all $x, y \in X$.}
\ans{(B)}

\subsection*{Q.21 -- Q.27 Multiple Select Question (MSQ), carry ONE mark each (no negative marking).}

\textbf{Q.21 (MSQ)} Let $P$ be a $3 \times 3$ real symmetric positive definite matrix. Which of the following statements is/are TRUE?
\optslistA{$Q^\top P Q$ is positive definite for all nonzero $Q \in \mathbb{R}^{3 \times 3}$.}{$Q^\top P Q$ is positive semidefinite for all $Q \in \mathbb{R}^{3 \times 3}$.}{$Q^\top P Q$ is positive definite if $Q \in \mathbb{R}^{3 \times 3}$ is nonsingular.}{$Q^\top P Q$ is not positive semidefinite for some $Q \in \mathbb{R}^{3 \times 3}$.}
\ans{(B), (C) (MSQ)}

\textbf{Q.22 (MSQ)} Let $f : \mathbb{R}^2 \to \mathbb{R}$ be a function defined as follows:
\[
  f(x, y) = \begin{cases} \dfrac{x\sqrt{x^2 + y^2}}{|x|}, & \text{for } x \ne 0, \\ 0, & \text{for } x = 0. \end{cases}
\]
Which of the following statements is/are TRUE?
\optslistA{$f(x, y)$ is continuous at $(0, 0)$.}{For any $u \in \mathbb{R}^2$, the directional derivative of $f$ at $(0, 0)$ in the direction of $u$ exists.}{$\dfrac{\partial f}{\partial x}(0, 0) = 1 = \dfrac{\partial f}{\partial y}(0, 0)$.}{$f(x, y)$ is differentiable at $(0, 0)$.}
\ans{(A), (B) (MSQ)}

\textbf{Q.23 (MSQ)} Let $\mathbb{Z}_{11}$ be the ring of integers modulo $11$. Let $\mathbb{Z}_{11}[x]$ be the ring of all polynomials with coefficients in $\mathbb{Z}_{11}$. Which of the following statements is/are TRUE?
\optslistA{The polynomial $x^2 + x + 4$ is irreducible over $\mathbb{Z}_{11}$.}{$\mathbb{Z}_{11}[x]/\langle x^2 + x + 4 \rangle$ is a field.}{The ideal $\langle x^2 + x + 4 \rangle$ is maximal.}{The ideal $\langle x^2 + x + 4 \rangle$ is not a prime ideal.}
\ans{(A), (B), (C) (MSQ)}

\textbf{Q.24 (MSQ)} Consider $\mathbb{R}$ with the topology $\tau = \{A \subseteq \mathbb{R} : A^c \text{ is finite}\} \cup \{\emptyset\}$. Which of the following statements is/are TRUE?
\optslistA{$(\mathbb{R}, \tau)$ is a Hausdorff space.}{Any finite subset of $(\mathbb{R}, \tau)$ is closed.}{$(\mathbb{R}, \tau)$ is compact.}{$(\mathbb{R}, \tau)$ is connected.}
\ans{(B), (C), (D) (MSQ)}

\textbf{Q.25 (MSQ)} Consider the following partial differential equation:
\[
  a \frac{\partial^2 f(x, y)}{\partial x^2} + b \frac{\partial^2 f(x, y)}{\partial y^2} = 8 f(x, y),
\]
where $a$ and $b$ are distinct positive real numbers. The combination(s) of the values of the real parameters $\xi$ and $\eta$ for which $f(x, y) = e^{2\xi x + \eta y}$ is a solution of the given partial differential equation, is/are
\optsfourA{$\xi = \dfrac{1}{\sqrt{a}}, \; \eta = \dfrac{2}{\sqrt{b}}$}{$\xi = 0, \; \eta = 0$}{$\xi = \dfrac{1}{\sqrt{2a}}, \; \eta = \dfrac{1}{\sqrt{2b}}$}{$\xi = 0, \; \eta = \dfrac{2\sqrt{2}}{\sqrt{b}}$}
\ans{(A), (D) (MSQ)}

\textbf{Q.26 (MSQ)} Consider the following assignment problem where $X, Y, Z$ are tasks, $P, Q, R$ are agents and the cost matrix is given by
\begin{center}
\begin{tabular}{cccc}
\toprule
& $X$ & $Y$ & $Z$ \\
\midrule
$P$ & $4$ & $2$ & $8$ \\
$Q$ & $2$ & $3$ & $7$ \\
$R$ & $3$ & $1$ & $6$ \\
\bottomrule
\end{tabular}
\end{center}
Which of the following statements is/are TRUE for an optimal assignment?
\optsfourA{The cost is $9$.}{Agent $P$ is assigned to task $Y$.}{Agent $Q$ is assigned to task $X$.}{Agent $R$ is assigned to task $Y$.}
\ans{(B), (C) (MSQ)}

\textbf{Q.27 (MSQ)} Consider the following differential equation:
\[
  2x(x - 2)^2 \frac{d^2 y}{dx^2} + 3x \frac{dy}{dx} + (x - 2)y = 0.
\]
Which of the following statements is/are TRUE?
\optslistA{The point $x = 2$ is not a singular point.}{The point $x = 2$ is a singular point.}{The point $x = 2$ is a regular singular point.}{The point $x = 2$ is not a regular singular point.}
\ans{(B), (D) (MSQ)}

\subsection*{Q.28 -- Q.35 Numerical Answer Type (NAT), carry ONE mark each (no negative marking).}

\textbf{Q.28} Let $T : \mathbb{R}^3 \to \mathbb{R}^3$ be the linear transformation which reflects every vector in $\mathbb{R}^3$ through a two-dimensional subspace of $\mathbb{R}^3$. Let $P \in \mathbb{R}^{3 \times 3}$ be the matrix representation of $T$ using the standard basis $\left\{\begin{pmatrix} 1 \\ 0 \\ 0 \end{pmatrix}, \begin{pmatrix} 0 \\ 1 \\ 0 \end{pmatrix}, \begin{pmatrix} 0 \\ 0 \\ 1 \end{pmatrix}\right\}$. Then the value of $2 \times \operatorname{trace}(P) - 3 \times \operatorname{determinant}(P)$ is equal to \underline{\hspace{2cm}} (answer in integer).
\ans{$5$}

\textbf{Q.29} Let $x^3 - x + 1 \in \mathbb{Z}_3[x]$, where $\mathbb{Z}_3[x]$ is the ring of all polynomials with coefficients in $\mathbb{Z}_3$. Then the degree of the field extension $\mathbb{Z}_3[x]/\langle x^3 - x + 1 \rangle$ of $\mathbb{Z}_3$ is equal to \underline{\hspace{2cm}} (answer in integer).
\ans{$3$}

\textbf{Q.30} Consider the differential equation $\dfrac{dy}{dx} = x + y$ with the initial condition $y(0) = 1$. Using the modified Euler's method, the second approximation to $y(h)$, where $h = 0.05$ (step size), is equal to \underline{\hspace{2cm}} (rounded off to TWO decimal places).
\ans{$1.00 \text{ to } 1.10$}

\textbf{Q.31} Let
\[
  \alpha = \lim_{n \to \infty} \int_0^1 \frac{n^2 + (\sin e^x)^n}{7n^2 + x^8}\,dx.
\]
The value of $14\alpha$ is equal to \underline{\hspace{2cm}} (answer in integer).
\ans{$2$}

\textbf{Q.32} Let $X = (\mathbb{R}^3, \|\cdot\|_1)$ where $\left\|\begin{pmatrix} x \\ y \\ z \end{pmatrix}\right\|_1 = |x| + |y| + |z|$ and $T : X \to X$ be a linear transformation defined as
\[
  T\begin{pmatrix} x \\ y \\ z \end{pmatrix} = \begin{pmatrix} 1 & 1 & 3 \\ 2 & 2 & -2 \\ 1 & 3 & -3 \end{pmatrix} \begin{pmatrix} x \\ y \\ z \end{pmatrix}.
\]
Then the operator norm of $T$ is equal to \underline{\hspace{2cm}} (answer in integer).
\ans{$8$}

\textbf{Q.33} Consider the following Linear Programming Problem (LPP):
\begin{align*}
  \text{maximize} \quad & z = 10x_1 + 20x_2 \\
  \text{subject to} \quad & x_1 \le 36, \quad x_2 \ge 42, \quad x_1 + x_2 \ge 48, \quad 5x_1 + x_2 \le 150, \quad x_1, x_2 \ge 0.
\end{align*}
The maximum value of $z$ in the above LPP is equal to \underline{\hspace{2cm}} (answer in integer).
\ans{$3000$}

\textbf{Q.34} The number of distinct topologies on the set $\{1, 2, 3\}$ consisting of exactly four elements is equal to \underline{\hspace{2cm}} (answer in integer).
\ans{$9$}

\textbf{Q.35} For $a_n \in \mathbb{C}$, $n = 0, 1, 2, \ldots$, the power series $\sum_{n=0}^{\infty} a_n (z - 2)^n$ converges at $z = 5$ and diverges at $z = -1$. Then the radius of convergence of this power series is equal to \underline{\hspace{2cm}} (answer in integer).
\ans{$3$}

\subsection*{Q.36 -- Q.47 Multiple Choice Question (MCQ), carry TWO marks each (for each wrong answer: $-2/3$).}

\textbf{Q.36} The solution $y(x) = \begin{pmatrix} y_1(x) \\ y_2(x) \end{pmatrix}$ of the initial value problem
\[
  \frac{dy}{dx} = \begin{pmatrix} -3 & 4 \\ -2 & 3 \end{pmatrix} y, \qquad y(0) = \begin{pmatrix} 1 \\ 2 \end{pmatrix}
\]
is equal to
\optsfourA{$3\begin{pmatrix} 1 \\ 1 \end{pmatrix} e^x + \begin{pmatrix} 2 \\ 1 \end{pmatrix} e^{-x}$}{$-3\begin{pmatrix} 1 \\ 1 \end{pmatrix} e^x + \begin{pmatrix} 2 \\ 1 \end{pmatrix} e^{-x}$}{$3\begin{pmatrix} 1 \\ 1 \end{pmatrix} e^x + 2\begin{pmatrix} 2 \\ 1 \end{pmatrix} e^{-x}$}{$3\begin{pmatrix} 1 \\ 1 \end{pmatrix} e^x - \begin{pmatrix} 2 \\ 1 \end{pmatrix} e^{-x}$}
\ans{(D)}

\textbf{Q.37} Let $P = \begin{pmatrix} 1 & 0 & 0 \\ 0 & 9 & 0 \\ 0 & 0 & 16 \end{pmatrix}$. Define an inner product on $\mathbb{R}^3$ with respect to $P$ as $\langle x, y \rangle_P = x^\top P y$, for all $x, y \in \mathbb{R}^3$. Consider the subspace $V = \operatorname{span}\left\{\begin{pmatrix} 1 \\ 0 \\ 0 \end{pmatrix}, \begin{pmatrix} 1 \\ 1 \\ 1 \end{pmatrix}\right\}$. Which one of the following sets is an orthonormal basis of $V$ with respect to $\langle x, y \rangle_P$?
\optsfourA{$\left\{\begin{pmatrix} 1 \\ 0 \\ 0 \end{pmatrix}, \begin{pmatrix} 0 \\ 1 \\ 1 \end{pmatrix}\right\}$}{$\left\{\begin{pmatrix} 0 \\ 1 \\ 0 \end{pmatrix}, \begin{pmatrix} 0 \\ 1 \\ 1 \end{pmatrix}\right\}$}{$\left\{\begin{pmatrix} 1 \\ 0 \\ 0 \end{pmatrix}, \begin{pmatrix} 0 \\ 1/5 \\ 1/5 \end{pmatrix}\right\}$}{$\left\{\begin{pmatrix} 1 \\ 0 \\ 0 \end{pmatrix}, \begin{pmatrix} 0 \\ 1/\sqrt{2} \\ 1/\sqrt{2} \end{pmatrix}\right\}$}
\ans{(C)}

\textbf{Q.38} The value of the contour integral
\[
  \int_\Gamma \frac{(1 - z)(3\cos z + 5\sin z)}{1 - z^{101}}\,dz,
\]
where $\Gamma = \left\{z \in \mathbb{C} : |z| = \frac{7}{9}\right\}$, oriented in the counter-clockwise direction, is equal to
\optsfourA{$2\pi i$}{$-2\pi i$}{$0$}{$4\pi i$}
\ans{(C)}

\textbf{Q.39} Let $f(z) = \sum_{n=1}^{\infty} 5^{-n} \cos(nz)$. On which one of the following domains, does $f(z)$ represent an analytic function?
\optsfourA{$\{z \in \mathbb{C} : |\operatorname{Im} z| < \ln 5\}$}{$\{z \in \mathbb{C} : |\operatorname{Im} z| > \ln 5\}$}{$\{z \in \mathbb{C} : |\operatorname{Re} z| < \ln 5\}$}{$\{z \in \mathbb{C} : |\operatorname{Re} z| > \ln 5\}$}
\ans{(A)}

\textbf{Q.40} Let $\delta(t)$ be the unit impulse function defined by $\delta(x - x_0) = 0$ when $x \ne x_0$ and $\int_{-\infty}^{\infty} \delta(x - x_0)\,dx = 1$. Consider the following initial value problem
\[
  2\frac{d^2 y}{dx^2} + \frac{dy}{dx} + 2y = \delta(x - 5)
\]
with $y(0) = 0$ and $\frac{dy}{dx}(0) = 0$. Which one of the following is TRUE?
\optsfourA{$y(10) = \dfrac{10}{\sqrt{15}} e^{-5/4} \sin\!\left(\dfrac{\sqrt{15}}{4}\right)$}{$y(10) = \dfrac{2}{\sqrt{15}} e^{-5/4} \sin\!\left(\dfrac{5\sqrt{15}}{4}\right)$}{$y(10) = \dfrac{10}{\sqrt{15}} e^{5/4} \sin\!\left(\dfrac{\sqrt{15}}{4}\right)$}{$y(10) = \dfrac{2}{\sqrt{15}} e^{5/4} \sin\!\left(\dfrac{5\sqrt{15}}{4}\right)$}
\ans{(B)}

\textbf{Q.41} Which one of the following statements is TRUE?
\optslistA{The set of all $2 \times 2$ real matrices under usual matrix addition and multiplication is a Unique Factorization Domain.}{The set of all real valued functions under usual function addition and multiplication is a Unique Factorization Domain.}{The set of all polynomials with coefficients in $\mathbb{R}$ under usual polynomial addition and multiplication is a Unique Factorization Domain.}{The set of Riemann integrable functions in $[0, 1]$ under usual function addition and multiplication is a Unique Factorization Domain.}
\ans{(C)}

\textbf{Q.42} Let $\ell_\infty = \{x = (x_n)_{n \ge 1} \mid x_n \in \mathbb{R}, \, \sup\{|x_n| : n = 1, 2, \ldots\} < \infty\}$ with the supremum norm. Let $T : \ell_\infty \to \ell_\infty$ be given by $T(x_1, x_2, x_3, \ldots) = \left(x_1, \frac{x_2}{2}, \frac{x_3}{3}, \ldots\right)$. Which one of the following is TRUE?
\optslistA{$T$ is bounded but not one-to-one.}{$T$ is one-to-one but not bounded.}{$T$ is bounded and the inverse (from the range of $T$) exists but not bounded.}{$T$ is bounded and the inverse (from the range of $T$) is bounded.}
\ans{(C)}

\textbf{Q.43} Let $u(x, t)$ satisfy the wave equation
\[
  \frac{\partial^2 u}{\partial t^2} = c^2 \frac{\partial^2 u}{\partial x^2}, \quad -\infty < x < \infty, \; t > 0, \; c > 0,
\]
with $u(x, 0) = \begin{cases} 1, & \text{if } |x| < 1, \\ 0, & \text{otherwise,} \end{cases}$ and $\dfrac{\partial u}{\partial t}(x, 0) = 0$. By using D'Alembert's formula, the maximum value of $u(0, t)$ for $t > 0$ is
\optsfourA{$1$}{$\dfrac{1}{2}$}{$\dfrac{1}{4}$}{$0$}
\ans{(A)}

\textbf{Q.44} Consider the Laplace equation $\dfrac{\partial^2 T}{\partial x^2} + \dfrac{\partial^2 T}{\partial y^2} = 0$, $0 < x < 1$, $0 < y < 1$, with the boundary conditions
\[
  T(x, 0) = x, \quad T(0, y) = y, \quad T(x, 1) = 1 + x, \quad T(1, y) = 1 + y.
\]
Then the value of $T\!\left(\frac{1}{2}, \frac{1}{3}\right)$ is equal to
\optsfourA{$\dfrac{7}{6}$}{$\dfrac{5}{6}$}{$\dfrac{1}{6}$}{$\dfrac{1}{2}$}
\ans{(B)}

\textbf{Q.45} Let $E_1$ and $E_2$ be subsets of a normed linear space $X$, and
\[
  E_1 + E_2 = \{x + y : x \in E_1, \, y \in E_2\}, \qquad E_1 \times E_2 = \{(x, y) : x \in E_1, \, y \in E_2\}.
\]
Which one of the following is NOT TRUE?
\optslistA{If either of $E_1$ or $E_2$ is open, then $E_1 + E_2$ is open.}{If $E_1$ and $E_2$ are convex, then $E_1 + E_2$ is convex.}{If $E_1$ and $E_2$ are closed, then $E_1 + E_2$ is closed.}{If $E_1$ and $E_2$ are connected, then $E_1 \times E_2$ is connected.}
\ans{(C)}

\textbf{Q.46} For a transportation problem, let $c_{ij}$ denote the unit cost of the cell $(i, j)$. Assume that $c_{11} = 10$, $c_{12} = 12$, $c_{21} = 11$ and $c_{22} = 13$. Let $\alpha_i$ and $\beta_j$, $i, j = 1, 2, 3$, represent the simplex multipliers associated with any basis corresponding to the unit cost $c_{ij}$. Assume that $\alpha_1 = x$, $\alpha_2 = x + 1$, $\beta_1 = y$ and $\beta_2 = y + 2$. The relative cost coefficient $d_{ij}$ is the difference between the current solution and the new improved solution. If $x = 4$, $c_{13} = 19$ and $\beta_3 = y + 5$, then which one of the following is TRUE?
\optsfourA{$y = 6$ and $d_{13} = 7$}{$y = 6$ and $d_{13} = 4$}{$y = 5$ and $d_{13} = 7$}{$y = 5$ and $d_{13} = 4$}
\ans{(B)}

\textbf{Q.47} Let $X$ and $Y$ be topological spaces and $f : X \to Y$ be a continuous and bijective mapping. Which one of the following statements is TRUE?
\optslistA{$f$ is a homeomorphism if $X$ and $Y$ are compact.}{$f$ is a homeomorphism if $X$ is Hausdorff and $Y$ is compact.}{$f$ is a homeomorphism if $X$ is compact and $Y$ is Hausdorff.}{$f$ is a homeomorphism if $X$ and $Y$ are Hausdorff.}
\ans{(C)}

\subsection*{Q.48 -- Q.54 Multiple Select Question (MSQ), carry TWO marks each (no negative marking).}

\textbf{Q.48 (MSQ)} Let $P$ and $Q$ be $3 \times 3$ nonzero real matrices. Assume that there exists a $3 \times 3$ real nonsingular matrix $S$ such that $S^{-1}PS$ and $S^{-1}QS$ are both upper triangular. Which of the following statements is/are TRUE?
\optslistA{The matrix $PQ - QP$ is invertible.}{The matrix $PQ - QP$ is nilpotent.}{The matrix $I + PQ - QP$ is invertible, where $I$ is the $3 \times 3$ identity matrix.}{If $PQ - QP$ is a nonzero matrix, then $PQ - QP$ is diagonalizable.}
\ans{(B), (C) (MSQ)}

\textbf{Q.49 (MSQ)} Let $S = \bigcup_{n=1}^{\infty} \left\{(x, y) \in \mathbb{R}^2 : (x - n)^2 + y^2 = \frac{1}{n^2}\right\}$ with usual topology. Which of the following statements is/are TRUE?
\optsfourA{$S$ is not compact.}{$S$ is compact.}{$S$ is connected.}{$S$ is not connected.}
\ans{(A), (D) (MSQ)}

\textbf{Q.50 (MSQ)} Let $f(z) = \dfrac{z}{1 - z}$ and $g(z) = \dfrac{1 + z}{1 - z}$ be two M\"obius transformations defined on the unit disc $D = \{z \in \mathbb{C} : |z| < 1\}$. Consider the following statements:\\
$S_1$: $f(D) \subseteq g(D)$\\
$S_2$: $g(D) \subseteq f(D)$\\
Which of the following statements is/are CORRECT?
\optslistA{$S_1$ is true.}{$S_2$ is true.}{$S_2$ is true and $S_1$ is false.}{Neither $S_1$ is true nor $S_2$ is true.}
\ans{(B), (C) (MSQ)}

\textbf{Q.51 (MSQ)} Let $\alpha$ and $\beta$ be the roots of the indicial equation obtained in the method of finding Frobenius series solution to the differential equation:
\[
  2x^2 \frac{d^2 y}{dx^2} - x \frac{dy}{dx} + (1 - x^2)y = 0.
\]
Which of the following statements is/are TRUE?
\optsfourA{$\alpha \ne \beta$}{$\alpha - \beta \in \mathbb{Z}$}{$\alpha - \beta \notin \mathbb{Z}$}{$\alpha = \beta$}
\ans{(A), (C) (MSQ)}

\textbf{Q.52 (MSQ)} Let $\mathbb{Q}[x]$ be the ring of all polynomials with coefficients in $\mathbb{Q}$ under the usual polynomial addition and multiplication. Let $T : \mathbb{Q}[x] \to \mathbb{Q}[x]$ be defined by $T(p(x)) = p(x^2)$, for all $p(x) \in \mathbb{Q}[x]$. Which of the following statements is/are TRUE?
\optsfourA{$T$ is a ring homomorphism.}{$T$ is one-to-one.}{$T$ is onto.}{$T$ is a ring isomorphism.}
\ans{(A), (B) (MSQ)}

\textbf{Q.53 (MSQ)} Let $X$ be any normed linear space and $X'$ be the dual space of $X$. Which of the following statements is/are TRUE?
\optslistA{If $X$ is separable and $X'$ is non-separable, then $X$ must be reflexive.}{If $X$ is separable and $X'$ is non-separable, then $X$ cannot be reflexive.}{If $X$ is reflexive, then $X'$ is reflexive.}{If $X'$ is separable, then $X$ is separable.}
\ans{(B), (C), (D) (MSQ)}

\textbf{Q.54 (MSQ)} Let $(\mathbb{Z}_n, +)$ be the group of integers modulo $n$. Let $G = \mathbb{Z}_{30} \oplus \mathbb{Z}_{12}$ be the external direct product of $\mathbb{Z}_{30}$ and $\mathbb{Z}_{12}$. Then which of the following statements is/are TRUE?
\optslistA{$G$ is not a cyclic group.}{The order of the element $(14, 7)$ in $G$ is $98$.}{The order of the element $(14, 7)$ in $G$ is $60$.}{There is an element in $G$ of order $360$.}
\ans{(A), (C) (MSQ)}

\subsection*{Q.55 -- Q.65 Numerical Answer Type (NAT), carry TWO marks each (no negative marking).}

\textbf{Q.55} Let
\[
  J = \begin{pmatrix}
    2 & 1 & 0 & 0 & 0 & 0 \\
    0 & 2 & 0 & 0 & 0 & 0 \\
    0 & 0 & 2 & 1 & 0 & 0 \\
    0 & 0 & 0 & 2 & 0 & 0 \\
    0 & 0 & 0 & 0 & 3 & 1 \\
    0 & 0 & 0 & 0 & 0 & 3
  \end{pmatrix}.
\]
Then the geometric multiplicity of the eigenvalue $2$ of $J$ is equal to \underline{\hspace{2cm}} (answer in integer).
\ans{$2$}

\textbf{Q.56} Let $\alpha = \iint_S \mathbf{F} \cdot \hat{\mathbf{n}}\,dS$, where $\mathbf{F} = (2x + 3z)\hat{\mathbf{i}} + (xz - y)\hat{\mathbf{j}} + (y^2 + 2z)\hat{\mathbf{k}}$ and $S$ is the sphere with centre at $(3, -1, 2)$ and radius $9$. Here, $\hat{\mathbf{n}}$ is the unit normal drawn outward and $\hat{\mathbf{i}}, \hat{\mathbf{j}}, \hat{\mathbf{k}}$ are unit vectors. Then the value of $\dfrac{1}{36\pi}\alpha$ is equal to \underline{\hspace{2cm}} (answer in integer).
\ans{$81$}

\textbf{Q.57} Consider the problem of maximizing
\[
  z = \begin{pmatrix} x_1 & x_2 & x_3 \end{pmatrix} \begin{pmatrix} 2 & 0 & 0 \\ 0 & 2 & 1 \\ 0 & 1 & 2 \end{pmatrix} \begin{pmatrix} x_1 \\ x_2 \\ x_3 \end{pmatrix}
\]
subject to $\begin{pmatrix} x_1 & x_2 & x_3 \end{pmatrix} \begin{pmatrix} x_1 \\ x_2 \\ x_3 \end{pmatrix} = 1$, where $\begin{pmatrix} x_1 \\ x_2 \\ x_3 \end{pmatrix} \in \mathbb{R}^3$. Then the maximum value of $z$ is \underline{\hspace{2cm}} (answer in integer).
\ans{$3$}

\textbf{Q.58} Let $P_3(\mathbb{R})$ be the vector space of all polynomials of degree at most three with the real coefficients under usual polynomial addition and scalar multiplication. Let $T : P_3(\mathbb{R}) \to \mathbb{R}^2$ be the linear transformation defined as $T(p) = (p(1), p'(1))$ for all $p \in P_3(\mathbb{R})$, where $p'$ is the derivative of $p$. Then the nullity of $T$ is equal to \underline{\hspace{2cm}} (answer in integer).
\ans{$2$}

\textbf{Q.59} The number of zeros of the complex polynomial $z^6 + 5z^3 + 4z + 11$ in the annulus $\{z \in \mathbb{C} : 1 < |z| < 3\}$ is equal to \underline{\hspace{2cm}} (answer in integer).
\ans{$6$}

\textbf{Q.60} Let $L^2[0, \pi]$ denote the space of all real valued Lebesgue square integrable functions on $[0, \pi]$. Let $T : L^2[0, \pi] \to L^2[0, \pi]$ be defined as follows:
\[
  T(f(x)) = \sin x \int_0^\pi f(t)\cos t\,dt + \cos x \int_0^\pi f(t)\sin t\,dt.
\]
Then the value of $\dfrac{4}{\pi}\|T\|$ is equal to \underline{\hspace{2cm}} (answer in integer).
\ans{$2$}

\textbf{Q.61} Let $\Omega = \{(x, y) \in \mathbb{R}^2 : x^2 + y^2 < 1\}$ be the open unit disc and $\partial\Omega$ be its boundary. If $u(x, y)$ is the solution of the following Dirichlet problem
\[
  \frac{\partial^2 u}{\partial x^2} + \frac{\partial^2 u}{\partial y^2} = 0 \quad \text{in } \Omega, \qquad u(x, y) = x^2 - y^2 \quad \text{on } \partial\Omega,
\]
then the value of $4\!\left(u\!\left(\frac{1}{2}, 0\right) - u\!\left(0, \frac{1}{2}\right)\right)$ is \underline{\hspace{2cm}} (answer in integer).
\ans{$2$}

\textbf{Q.62} Let $X = \{5, 6, 7, 8, 9, 10\}$ be equipped with the topology $\tau = \{\emptyset, X, \{5, 6, 7\}, \{8, 9, 10\}\}$. Then the number of subsets of $X$ which are neither open nor closed is \underline{\hspace{2cm}} (answer in integer).
\ans{$60$}

\textbf{Q.63} Let $\alpha, \beta \in \mathbb{R}$. If $(4, 0, 2, \beta)$ is an optimal solution of the Linear Programming Problem:
\begin{align*}
  \text{minimize} \quad & x_1 + 3x_2 + 2x_3 - \alpha x_4 \\
  \text{subject to} \quad & 4x_1 + x_2 + x_3 = 18, \\
  & -3x_1 + 2x_3 + x_4 = 2, \\
  & x_1, x_2, x_3, x_4 \ge 0,
\end{align*}
then the maximum value of $22(\alpha + \beta)$ is equal to \underline{\hspace{2cm}} (answer in integer).
\ans{$206$}

\textbf{Q.64} Let $D = \{(x, y) \in \mathbb{R}^2 : 0 \le x \le 2, \; 0 \le y \le 2\}$ and let $f(t)$ denote the smallest integer greater than or equal to $t$. Then the value of the integral $\iint_D f(x + y)\,dx\,dy$ is \underline{\hspace{2cm}} (answer in integer).
\ans{$10$}

\textbf{Q.65} If Jacobi method is used to solve the following system of linear equations
\[
  \begin{pmatrix} 1 & 2 & 1 \\ 0 & 2 & 2 \\ 1 & 1 & 1 \end{pmatrix} \begin{pmatrix} x_1 \\ x_2 \\ x_3 \end{pmatrix} = \begin{pmatrix} 2 \\ 2 \\ 2 \end{pmatrix}
\]
with the initial guess $x^{(0)} = \begin{pmatrix} 0 \\ 0 \\ 0 \end{pmatrix}$ and $x^{(i)} = \begin{pmatrix} x_1^{(i)} \\ x_2^{(i)} \\ x_3^{(i)} \end{pmatrix}$, $i = 1, 2, 3, \ldots$, denotes the $i$th iterate, then the value of $|x_1^{(2)} + x_2^{(2)} + x_3^{(2)}|$ is equal to \underline{\hspace{2cm}} (answer in integer).
\ans{$4$}

\end{document}
